\section{CM period and determinant comparison}\label{sec:cm}

This section proves the CM comparison, Proposition~\ref{prop:cm}, and the
paired CM determinant formula needed for the imaginary \(S_3\) anchor.
There are two distinct rank arguments. First, under a \emph{known} simple
analytic zero, the Heegner argument \eqref{eq:K0} supplies the Selmer line
and permits an exact period and logarithm calculation. Second, starting
only with \(s_2(E)=1\), a tame Bockstein test detects a simple zero.
The paired formula belongs to the first argument and is independent of
the second.

\subsection{CM realizations and period normalization}

Let \(L\) be an imaginary quadratic field and let \(f\) be a primitive
weight-two CM newform induced from \(L\), with trivial nebentypus and
level \(N\). Write \(F\) for its Hecke field and \(F_c=FL\).
We use either of the following realizations:
\begin{enumerate}[label=(\roman*)]
\item \(A=e_fJ_0(N)\), with \(e_f\) the rational orbit projector,
the polarization \(\lambda_A\) restricted from \(J_0(N)\), and the
algebraic eigendifferential \(\omega\) pulling back to \(f\,dq/q\);
\item a CM elliptic curve \(A/\mathbb Q\), with its principal
polarization, its primitive newform \(f\), \(F=\mathbb Q\), and a
minimal N\'eron differential \(\omega\).
\end{enumerate}
In the first case the projector and maps to the orbit factor are taken
up to isogeny. Fix compatible complex and dyadic embeddings of the
algebraic data. After enlarging the dyadic coefficient field when
necessary, write \(\mathcal O\) for its ring of integers and \(\pi\)
for a uniformizer.

Rosati fixes \(F\). The polarization pairings, their rational inverses
on dual Tate modules, and the Poincar\'e height pairing therefore have
\(F\)-valued versions whose traces to \(\mathbb Q\) are the ordinary
pairings. At a single coefficient embedding one computes the pairing
by projecting one argument and applying the ordinary pairing.

\begin{lemma}[CM realizations]\label{cm:realization}
The field \(F\) is real, the Hecke-character value field is \(F_c\),
and the characteristic-zero Tate plane at each Hecke embedding is
absolutely irreducible over \(\mathbb Q\).
The variety \(A\) is \(\mathbb Q\)-simple, with endomorphism algebra
\(F\); over \(L\) its endomorphism algebra is \(F_c\), and its CM type
is induced from \(L\).
\end{lemma}

\begin{proof}
By classical CM and theta-series theory
\cite[Theorems~1.1 and~1.4]{Schuett09}, \(N\) is divisible by the primes
ramified in \(L\);
the Hecke character \(\varphi\) has type \((1,0)\) in the convention \(\varphi((x))=x\) near 1
modulo the conductor (exchange orientations if necessary). Its conjugate-source character
equals its complex conjugate, by purity and the trivial determinant character of the form: on
rational arguments \(n\) away from the modulus, \(\varphi((n))=n\chi_L(n)\) for positive
\(n\). Thus \(F\) is real.

The field of values of \(\varphi\) on good ideals is \(F_c\). Indeed an automorphism fixing
\(F,L\) changes it at most by a finite character, and leaves split-prime traces and products
unchanged. Interchanging the two values at a split prime cannot give a root-of-unity ratio
(power up to ray principal ideals). Thus the finite character is trivial by Chebotarev, and
the reverse generation follows by traces and ray-principal values.

The Tate realizations are inductions of the character lines by the CM theorem
\cite[Sections~3--4]{Ribet77}, or by traces at good primes and semisimplicity.
On each Hecke embedding the representation over \(\mathbb Q\)
is absolutely irreducible and the two lines over \(L\) differ to infinite order. For distinct
embeddings the planes are inequivalent, so also no lines between them can coincide over \(L\).

Faltings' theorem \cite{Faltings83} implies that \(A\) is \(\mathbb Q\)-simple with
endomorphism algebra \(F\), and over \(L\) its endomorphisms have algebra \(F_c\). Indeed this
latter algebra contains \(F\); locally at all finite coefficient places its base changes are
the commutants of the separate character pairs, the quadratic étale algebras of \(F_c/F\) by
the character-value description. This determines the quadratic étale algebra globally. The
\(F\)-Lie algebra over \(\mathbb Q\) has rank one, so the action of the full field on it over
\(L\) induces CM type from \(L\).
\end{proof}

Write \(V=H^1_{\rm B}(A)\) and orient labels so that
\[
    V\otimes_F F_c=V_\alpha\oplus c V_\alpha,\quad
    F_c\gamma=V_\alpha,\quad [\omega]_{\rm per}=p_\gamma\,\gamma .
\]
Here \(c\) is the real involution acting linearly (and its étale transport from a chosen
complex conjugation), \(\gamma\) a coefficient basis of one CM eigensummand, and the last
equality uses the chosen complex component and ordinary integration into Betti cohomology. In
compatible étale notation use
\[
       T_\alpha=\mathcal O\gamma(1),\qquad
       M^*=T_\alpha\oplus cT_\alpha
              =\operatorname{Ind}_{G_L}^{G_{\mathbb Q}}T_\alpha .
\]
The rationalized \(M^*\) is dual abelian Tate (\(T_2(A^\vee)\) with projection and enlarged
scalars). Put \(\mu=\langle\gamma,c\gamma\rangle_F\) by the inverse polarization Betti
pairing; in the projected Jacobian case this uses the cup of the two pullbacks to the curve.
Indeed the projector is orthogonal, so the inverse pairing for the restricted polarization is
just the pullback pairing via that projector. The corresponding dual-Tate alternating
polarization form divided by \(\mu\) is unimodular on \(M^*\) (in compatible twist bases,
ignoring sign). In the projected newform normalization we have
\[
                   \mu p_\gamma^2=-i\,8\pi^2(f,f)_N,                 \tag{C1}\label{eq:C1}
\]
with cup sign chosen accordingly: \(c[\omega]=[\bar\omega]\) on this component, and the
integral of the wedge computes exactly the stated Petersson pairing, with no Hecke-field trace
multiplicity.

\begin{lemma}[A CM lift of an imaginary residual plane]\label{cm:residual-lift}
If an \(S_3\) residual plane \(W/\mathbb F_2\) has imaginary quadratic
subfield \(L\), there is a form \(f\) as above for which
\(M^*/\pi\simeq W\otimes_{\mathbb F_2}\mathcal O/\pi\).
\end{lemma}

\begin{proof}
Start with a large rational conjugation-invariant modulus containing the discriminant;
prescribe a multiplicative residue character \(\eta(n)=\chi_L(n)\) on rational units and
\(\eta(e)=e^{-1}\) on global units. For sufficiently divisible modulus the images intersect
only at \(\pm1\), compatibly. Extend to all residues that are units. Prescribe \(x\eta(x)\) on
the group of prime-to-modulus principal ideals \((x)\); extend to prime-to-modulus fractional
ideals by extracting roots. This gives type \((1,0)\) with trivial determinant character on
the theta form (at its primitive level), cuspidal by non-invariance.

The residual inducing character of the plane has odd order and is inverted by conjugation.
Multiply the Grössencharacter by a finite odd Teichmüller character to arrange that the new
residual inducing character equals the desired order-three one. This uses abelian class field
theory to twist by the quotient (or inverse according to realizations). The finite character
is anticyclotomic and is trivial on rational positive arguments off the modulus (split pairs
by inversion, inert primes by inversion and odd order). Thus it does not change the
determinant character. The two residual orientations give the same induction.
\end{proof}

\subsection{Elliptic-unit inputs}

We use the norm-compatible rational elliptic-unit system
\(z_{2^\infty\mathfrak f}\) in the ray towers of \(L\), with the
normalization in \cite[Sections~3--4]{BF23}.

\begin{theorem}[Individual-character elliptic-unit comparison]
\label{cm:external-units}
Let \(\Gamma\simeq\mathbb Z_2^2\) be the free pro-\(2\) factor of the ray-tower group,
let \(\Lambda_{\mathcal O}=\mathcal O[[\Gamma]]\), and let
\(\Lambda(\chi)(1)\) have the finite-character and regular variable
actions of \cite[Section~4.1]{BF23}. If
\(\mathfrak f_\chi\mid\mathfrak f\), then, inside generic
rank-one cohomology,
\[
 \Lambda_{\mathcal O}z_{2^\infty\mathfrak f}
 =
 \det_{\Lambda_{\mathcal O}}^{-1}
 R\Gamma\bigl(\mathcal O_L[1/2\mathfrak f],\Lambda(\chi)(1)\bigr).
\]
For the character motive of type \((1,0)\) in the ideal-character
convention of Lemma~\ref{cm:realization}, a modulus divisible by its
conductor, and a compatible Betti vector \(\gamma\), elliptic-unit
reciprocity gives the period equation
\[
 \operatorname{per}\bigl(\exp^*z(\gamma)\bigr)
       =L_{2\mathfrak f}(\bar\varphi,1)\gamma
\]
on the holomorphic component, with the corresponding labeling of
\(\varphi\). The equation includes the case of a zero \(L\)-value.
\end{theorem}

The determinant assertion is the Johnson-Leung--Kings theorem in the
form of \cite[Theorem~4.1]{BF23}, applied to one finite character at a
time. That theorem includes \(p=2\).
The period equation is Kato's reciprocity law
\cite[Proposition~15.9]{Kato04}, in the CM formulation
\cite[Proposition~3.1]{BF23}. The latter uses infinity type
\((-1,0)\) in its convention; here the ideal-character convention is
\(\varphi((x))=x\) for \(x\) sufficiently close to \(1\) modulo the
conductor. In an abelian-variety realization the period map is ordinary
integration into Betti cohomology, projected at the embedding extending
the chosen embedding of the CM-type field. No factor two or base-field
trace is inserted in this period map.

We also use the distribution relations of this elliptic-unit system.
Dropping a fresh good prime ideal \(Q\) inserts
\(1-\mathrm{Fr}_Q^{-1}\) for the arithmetic action on units.
Auxiliary-ideal operators give integral smoothed units; the smoothing
ideal is prime to \(6\) and to all moduli. Dyadic conductor powers may
be increased throughout. These are the Robert-unit/Siegel-function
distributions with the normalization of
\cite[Sections~3.2 and~4.1]{BF23}.

The modulus \(\mathfrak f\) always contains the CM conductor and
the additional allowed support. Character maps on units and on twisted
cohomology are unnormalized coefficient push-forwards.

\begin{lemma}[Specialization under concentration]\label{cm:specialization}
Let \(\psi\) be an additional finite allowed character, possibly
trivial. Suppose that
\[
 R\Gamma\bigl(\mathcal O_L[1/2\mathfrak f],T_\alpha\psi\bigr)[1/2]
\]
has a one-dimensional \(H^1\) and no other cohomology.
Then the unsmoothed elliptic-unit class satisfies
\[
 \mathcal O z(T_\alpha\psi)=
 \mathcal D\bigl(R\Gamma(\mathcal O_L[1/2\mathfrak f],
                                      T_\alpha\psi)\bigr).
 \tag{C2}\label{eq:C2}
\]
Here the determinant is identified with the rational cohomology line.
No nonvanishing hypothesis on a complex \(L\)-value is required.
\end{lemma}

\begin{proof}
Split the ray-tower group into its torsion and free factors. Choose
\(\Lambda(\chi)\) from the torsion character of \(T_\alpha\psi\)
before the Tate twist, and evaluate the free factor at the remaining
character. The conductor of \(\chi\) away from \(2\) is accounted for
by finite inertia. Increase the dyadic part of \(\mathfrak f\), by
cofinality, until the full divisibility
\(\mathfrak f_\chi\mid\mathfrak f\) holds. Use the vector \(\gamma\)
as basis, enlarging \(\mathcal O\) if necessary.

Write
\[
 C_\Lambda=R\Gamma\bigl(\mathcal O_L[1/2\mathfrak f],
                                     \Lambda(\chi)(1)\bigr).
\]
Let \(\mathfrak q\) be the characteristic-zero specialization prime
of \(\Lambda_{\mathcal O}\). The residue-field complex at
\(\mathfrak q\) is concentrated in one line in degree one by hypothesis.
Cancellation in a minimal perfect model therefore identifies the
localized complex itself with one free term in that degree.
Consequently the generic determinant generator of
Theorem~\ref{cm:external-units} specializes to its actual
cohomological push-forward. This comparison may be made first with
a smoothing multiplier nonzero at \(\mathfrak q\), and then with
the rational unsmoothed class.

The integral determinant line specializes by derived base change:
\[
 \mathcal D(C_{\Lambda})\otimes_{\Lambda_{\mathcal O}}\mathcal O
 \simeq
 \mathcal D\bigl(R\Gamma(\mathcal O_L[1/2\mathfrak f],
                                      T_\alpha\psi)\bigr).
\]
The specialization of the chosen elliptic-unit determinant generator
has just been identified, after inverting \(2\), with the distinguished
class \(z(T_\alpha\psi)\). Hence equality of these lattices is precisely
\eqref{eq:C2}. The arithmetic/étale comparison and finite models on
the imaginary base justify the displayed cohomological base change.
The argument uses concentration rather than the nonzero-\(L\)-value
hypothesis of \cite[Proposition~4.1]{BF23}.
\end{proof}

\begin{lemma}[Shapiro and the CM period equation]\label{cm:period-transfer}
Suppose \(\psi\) comes from a Dirichlet character on \(\mathbb Q\),
is locally trivial at \(2\), and the omissions are conjugation-stable.
Allow all ramification primes and write \(z\) also for the Shapiro
class in \(M^*\psi\) over \(\mathbb Q\).
After trivializing at \(2\) by a compatible geometric Artin basis,
\[
 \exp^*z=\frac{L_{2\mathfrak f}(f,\psi^\pm,1)}{p_\gamma}\,\omega .
 \tag{C3}\label{eq:C3}
\]
The sign \(\pm\) fixes the reciprocal convention on the finite
character; one may invert the variable throughout.
\end{lemma}

\begin{proof}
The notation for omissions here and later means omit the indicated prime ideals over \(L\)
from its character \(L\)-function (below whole sets over rational primes, giving omitted
factors of the primitive newform twist).

Apply the reciprocity equation of Theorem~\ref{cm:external-units} to the type \((1,0)\)
character modified by the finite twist, using \(\gamma\otimes e_\psi\) with \(e_\psi\) an
Artin basis normalized on the identity label, in compatible Betti and étale comparisons.
Indeed this motive over \(L\) occurs with coefficients on \(H^1(\operatorname{Res}_{H/L}
A_H)\), for \(H\) cutting out the finite character, projecting the CM and finite Artin
factors.

To compare to the CM realization in the reciprocity law one needs only rational isomorphisms:
the type \((1,0)\) character motive is realized on isogeny factors via Weil restriction of CM
elliptic curves (Hecke-character construction), and the projected Tate character lines agree.
Faltings' theorem \cite{Faltings83}, projected over \(L\), then supplies an algebraic map with
extended number-field coefficients, nonzero on this component (tensor at the chosen prime then
choose a nonzero map before completion), hence an isomorphism for the period comparisons.

The cohomology push of units twists exactly by the indicated vector, as tensor twisting and
taking images compose. Conjugating the Hecke character in the law gives the same modular
\(L\)-series, with the consistent finite norm-character twist. Under Shapiro, restriction and
projection to \(V_\alpha\) preserve \(\exp^*\); restriction compatibility uses the trace in
local duality, without a doubling. Only one base embedding over 2 contributes to the
\(H^0(\Omega)\) of this component, the one matching the holomorphic period, also when \(2\) is
nonsplit, after extending scalars.

An Artin de Rham basis, evaluated geometrically over the splitting field on this label, uses a
nonzero algebraic factor \(b\) times \(e_\psi\); its complex and dyadic comparisons there use
the two embeddings of the very same \(b\). It thus cancels when taking the scalar relative to
\(\omega\) by the stated local trivialization. This proves \eqref{eq:C3}; the coefficient
push-forward is unnormalized.

Similarly for a fixed quadratic character split at 2 one may think of \(B=A^\varepsilon\) with
transported polarization, basis, and \(\omega\); with these geometric identifications (extend
base scalars for \(\omega\) if needed) the denominator is again \(p_\gamma\).
\end{proof} 

\par\smallskip\noindent\textbf{Smoothing and Euler factors.}\enspace In moving sequences use a fixed nonzero smoothing,
denoted by \(\Delta\) at character evaluations, to give integral bounded cohomology, before
any finite traces. Its Euler-system multiplier can be kept constant and nonzero in
\eqref{eq:C3} by killing the norm of the smoothing ideal by the moving Dirichlet characters
(difference of the norm and the auxiliary translate multiplier, nonzero by CM purity). A fixed
denominator-clearing integer if used here is included in that same \(\Delta\).

The norm relation on dropping a good prime ideal \(Q\), after twisting by \(\gamma\), at total
splitting in the remaining finite tower uses the single factor
\(1-\rho_{T_\alpha}(\mathrm{Fr}_Q)/\mathrm NQ\): transporting \(\gamma\) across the
inverse-Frobenius translate inserts precisely its Frobenius scalar. Similarly with additional
finite twists. At a good split rational prime \(r\), dropping both ideals at trivial new
character thus gives \(1-a_r/r+1/r\) (for the actual form on the fixed branch).


\subsection{The horizontal logarithm at a known simple zero}

Let \(B=A\), or let \(B\) be a fixed quadratic twist of \(A\) split
at \(2\). Write \(f_B\), \(N_B\), and \(\lambda_B\) for its primitive
form, level, and polarization, and transport \(\omega\) in the twist.
Let \(S_f\) contain \(2\) and all necessary bad and ramified primes.
For the unsmoothed Shapiro class with these omissions write \(z_B\), and put
\[
 e_S=\prod_{q\in S_f}L_q(f_B,1)^{-1}.
\]

\begin{lemma}[Nonvanishing of the eigendifferential logarithm]
\label{cm:log-nonzero}
For every nontorsion \(P\in B(\mathbb Q)\), one has
\(\log_\omega P\ne0\).
\end{lemma}

\begin{proof}
The Tate characters and Faltings' theorem \cite{Faltings83} identify \(B\), up to
\(F\)-equivariant isogeny, with the orbit factor of \(f_B\).
Good-prime traces show that its Hecke field is still \(F\).
In particular \(B\) is \(\mathbb Q\)-simple.

Multiply \(P\) into the domain of the \(2\)-adic logarithm.
Its full logarithm is nonzero: the logarithm is injective on a
sufficiently small formal neighborhood, so a point in its kernel is
torsion. If \(\log_\omega P=0\), this nonzero full logarithm lies in
the proper algebraic subspace \(\ker\omega\) of the Lie algebra.
The \(p\)-adic analytic subgroup theorem
\cite[Theorem~2.2]{FuchsPham15} gives an algebraic subgroup \(H\)
containing this multiple of \(P\), with
\(\operatorname{Lie}H\subseteq\ker\omega\).
The Zariski closure of the subgroup generated by a rational
nontorsion point is defined over \(\mathbb Q\); by
\(\mathbb Q\)-simplicity it is all of \(B\).
Thus \(H=B\), contradicting the properness of \(\ker\omega\).
\end{proof}

\begin{proposition}[Horizontal CM logarithm]\label{cm:tame-log}
Suppose that \(L(f_B,s)\) has a simple zero at \(1\).
Then the rational compact Selmer group at the chosen coefficient
embedding is one-dimensional, generated by any nontorsion rational
point \(P\) after extension of scalars. Write
\(z_B=\alpha\lambda_{B,*}(P)\) in projected Kummer notation.
With the absolute \(F\)-valued Poincar\'e height pairing,
\[
 \alpha=\pm e_S\log_\omega P\,
       \frac{L'(f_B,1)}{p_\gamma H_{\lambda_B,F}(P,P)}.
 \tag{C4}\label{eq:C4}
\]
The height-normalized complex expression is interpreted dyadically
through the algebraic comparison in the proof.
\end{proposition}

\begin{proof}
Choose a split Heegner imaginary companion \(K'\) with nonzero
central value, as in \eqref{eq:T1}, at level \(N_B\).
It differs from \(L\): a prime ramified in \(L\) divides \(N_B\)
and must split in \(K'\). The induced plane remains absolutely
irreducible over \(K'\), because its two characters are still distinct
on the compositum. The Heegner argument \eqref{eq:K0}, applied to the
known simple product zero, gives the rational Selmer line.
The nonzero rational-sign Hilbert trace and Kummer injectivity give
\(\dim_F B(\mathbb Q)_{\mathbb Q}=1\).
By \eqref{eq:C3}, the class \(z_B\) belongs to that Selmer line.
Indeed, its dual exponential vanishes, and odd local \(H^1\)
vanishes by local Tate-torsion finiteness and duality.
Lemma~\ref{cm:log-nonzero} gives \(\log_\omega P\ne0\).

\par\smallskip\noindent\phantomsection\label{cm:construction:exponential}%
\textbf{Choice of primes and the exponential coordinate.}\enspace

Take the auxiliary index prime \(\ell\) there inert in \(K'\), split in \(L\), away from the
operators and bad primes. Such \(\ell\) has \(a_{f_B}(\ell)\ne0\) by the ray principal value
argument. Use primes \(r_i\) and even tame characters from \(\lambda_i:(\mathbb
Z/r_i)^\times\to\mathbb Z/2^{m_i}\), \(m_i\to\infty\), with the prescriptions of the
horizontal comparison, splitting also in a fixed finite normal \(F_0\) containing \(L\), the
twisting data, Hilbert class field of \(K'\), and fields imposing fixed congruences. Kill
fixed support, smoothing norms and the additional finite list of primes as there. Take
smoothed classes from \eqref{eq:C3} in the corresponding group coefficients (include the split
pair of ideals at \(r_i\)), with polynomial exponential coordinate \(Z_i\) orienting the
variable to interpolate \(\bar\chi\) for characters \(\chi=t^{\lambda_i}\). These classes and
the exponential polynomials have bounded denominators, by split group coefficients at 2. For
the primitive characters used in \eqref{eq:T4}–\eqref{eq:T5} they evaluate on the exponential
side as
\(\Delta e_S L(f_B,\bar\chi,1)/p_\gamma\), by \eqref{eq:C3}; at \(t=1\) the ordinary class is
\(\Delta e_r z_B\), \(e_r=1-a_{f_B}(r_i)/r_i+1/r_i\). In particular \(Z_i(1)=0\).
Write \(Z(v)\) for the resulting bounded-denominator coefficientwise limit of
\(Z_i(1+v)\), and \(Z'\) for its linear coefficient.

\par\smallskip\noindent\phantomsection\label{cm:construction:transversality}%
\textbf{Frobenius and height transversality.}\enspace We can prescribe \(r_i\to1\) and limiting
Frobenius \(g\) with no root-of-unity Tate eigenvalue on any dyadic component, determinant one
on each plane. In fact add to \(F_0\) all dyadic roots of unity and dyadic radicals of the
primes to be killed, obtaining \(D_0\). Its maximal abelian part over \(F_0\) has bounded
exponent over the cyclotomic tower, by conjugation on the translations by one fixed cyclotomic
scalar differing from one. Each CM character has infinite image over \(F_0\) on the cyclotomic
kernel (otherwise compare conjugation transports using normality and inversion on the infinite
anti-ratio). The image is therefore still infinite on \(G_{D_0}\), by abelianness over
\(F_0\).

Avoid the finitely many characters' torsion preimages simultaneously (closed with empty
interior). Frobenii approximating such \(g\in G_{D_0}\) are realized by Chebotarev, with the
roots and radicals split to sufficient simultaneous depth; use the same primitive root/residue
exponent convention as before. Write \(a_*\) for the limiting trace on the chosen plane.

We can additionally arrange a nonzero \(F\)-height \(h_{\lambda}^{\mathbb Q}(P,P)\) for this
tame limit character, with polarization \(\lambda_B\) and the specified coefficient projection
in this notation. Indeed form the biextension lifts used before with first argument \(P\)
(Kummer \(x\) in full Tate), taking linear combinations of dual arguments by rational
\(F\)-endomorphisms and scalar extension so as to project to our embedding (clear fixed
denominators).

The resulting block and fiber cochains \(\beta,k\) have \(dk=-\beta\cup x\), with \(\beta\)
representing the projected \(\lambda_{B,*}(P)\) up to sign; ordinary evaluations with \(x\)
compute the projected pairing. The projected \(x\) remains nonzero on \(G_{F_0(T_2B)}\), by
finite-extension restriction in characteristic zero and central homothety vanishing on the
image. A homothety acts simultaneously by a nonidentity scalar \(m\): multiply an element over
\(F_0\) of cyclotomic value \(m\ne1\) by its complex-conjugation transport. Adding the radical
tower cannot kill the restriction: its translation quotient has conjugation by this element of
homothety action according to the cyclotomic scalar \(m^2\), precluding such a nonzero
continuous equivariant map. The joint kernel is normal over \(\mathbb Q\); the projected
\(x\)-values there span the plane. The commutator of two elements there has zero \(x,\beta\),
and its \(k\)-value is up to sign twice the projected symplectic pairing of the \(x\)-values,
hence can be nonzero. This adjusts \(g\), preserving the preceding requirements, to give
\[
                k(g)-\beta(g) g_T (g_T-1)^{-1}x(g)\ne0
\]
where \(g_T\) is the action on Tate. The division calculation of \eqref{eq:T2}, valid on each
of the biextensions in the combination and with \((g_T-1)^{-1}\) bounded on all components at
the approximating primes, gives \(h_{\lambda}^{\mathbb Q}(P,P)\) equal to this value up to
sign. All retrace rules for these heights use just the two Poincaré laws; norm is unramified
over the same ring class fields of \(K'\), and pre-division downstairs uses the uniform 2-part
bound for the residue-group size of \(B\).

\par\smallskip\noindent\textbf{The cohomological derivative.}\enspace The cup computation modulo the squared parameter, as
in \eqref{eq:T2}, gives
\begin{equation}\label{cm:eq:cup-derivative}
 Z'\log_\omega P
   =\pm\Delta(2-a_*)\alpha h_{\lambda}^{\mathbb Q}(P,P).
\end{equation}
Indeed after a fixed multiplier the degree-one derivative term of the cocycle satisfies \(d
z_1=-\lambda_i\cup z_0\) up to character orientation with \(z_0=\pm\Delta e_r\alpha\beta\);
make the adjustment in full dual Tate with coefficients if needed, the cohomological equality
being by the norm and Selmer line and torsion bounded. Cup with \(x\) and add the
corresponding \(\lambda_i\cup k\). At 2 the cup uses ordinary differential evaluation
(projected class against \(x\)), and at \(r_i\) one trivializes \(x\) by the
bounded-denominator Frobenius divisions, giving exactly the fiber calculation; remaining local
terms have limit zero as before.

\par\smallskip\noindent\textbf{The spectral and intersection derivative.}\enspace Use the holomorphic and intersection
calculation of \eqref{eq:T3}--\eqref{eq:T6} at level \(N_B\), on the block of \(f_B\). It
yields
\begin{equation}\label{cm:eq:spectral-derivative}
 \begin{split}
 C_Bp_\gamma L(f_B\otimes\varepsilon_{K'},1)
                 \frac{Z'}{\Delta e_S}
       &=\pm(2-a_*)h_{\rm pr}^{K'}(P_0,P_0),\\
 C_B&=\frac{\sqrt{|\operatorname{disc}K'|}}
                     {8\pi^2(f_B,f_B)_{N_B}}.
 \end{split}
\end{equation}
Here \(P_0\) is the projected Hilbert trace on the Jacobian orbit
factor with its restricted polarization, and \(h_{\rm pr}\) is the
corresponding \(F\)-valued tame height at the chosen embedding.
The two comparison steps are as follows.
\begin{enumerate}[label=(\roman*)]

\item Use the isolating Hecke polynomial before \eqref{eq:T5} with algebraic coefficients,
clearing fixed denominators. The Gram and difference calculation is identical at the indicated
embedding. For the second \(L\)-factor use a Betti period \(p'\) for the primitive companion
form as in \eqref{eq:C1} and the same additive modular-symbol polynomials. It differs from
\(p_\gamma\) by a nonzero algebraic multiple by quadratic twisting, isogeny and the CM
eigen-differentials. Likewise \(p_\gamma\) is algebraically proportional to the Betti period
in \eqref{eq:C1} at level \(N_B\). Thus \(C_B p_\gamma p'\) is algebraic; the
period-normalized companion modular symbols have bounded denominators by the fixed relative
period lattices (Manin--Drinfeld and the CM Betti line). These facts justify the
bounded-coefficient interpolation and differentiation giving \eqref{eq:T5}, now with
\(p_\gamma,\Delta e_S\) in place of \(\Omega_0,d_0\).
\item The finite intersection matching is on the same modular curve and oriented Hilbert orbit
of \(K'\), before projection. Inserting the Hecke projector into the pairing applies the
restricted orbit polarization split by \(F\) (no divisor by a map degree in this notation).
One can project by linear combinations of the correspondences after the unramified pull-ups in
that proof. Dividing the scaled first argument downstairs on the orbit factor and retracing
over \(K'\) then works by the same biextension rules and the uniform residue torsion bound
from the Frobenius choices. This gives \eqref{cm:eq:spectral-derivative} after the
cancellation
in \eqref{eq:T4}.
\end{enumerate}
This spectral calculation uses the Hilbert trace itself and the
nonzero companion value. It does not use \(P\), the Selmer-line
identification, or the nonzero-height prescription.

\par\smallskip\noindent\textbf{Comparison with the absolute height.}\enspace In the simple-zero case, \eqref{eq:GZ} gives
\[
 C_B L'(f_B,1)L(f_B\otimes\varepsilon_{K'},1)
       =2H_{\rm pr}(P_0,P_0),
\]
with the absolute projected pairing. An \(F\)-equivariant isogeny to \(B\) sends \(P_0\) into \(F P\) rationally
(also by the sign and rank in \eqref{eq:K0}); its pullback of \(\lambda_B\) differs from the
restricted polarization by an element of \(F^\times\). Thus \(H_{\rm
pr}(P_0,P_0)/H_{\lambda_B,F}(P,P)\) is nonzero algebraic, and \(h_{\rm pr}^{K'}(P_0,P_0)\) is
twice that multiple of \(h_\lambda^\mathbb Q(P,P)\), by bilinearity, adjunction and extension
from \(\mathbb Q\) to \(K'\). Consequently the complex height-normalized \(L'/p_\gamma\) in
\eqref{eq:C4} is interpreted dyadically by an algebraic scalar, and the two derivative
computations imply \eqref{eq:C4}, cancelling the nonzero factors. This uses \eqref{eq:GZ} with
its absolute-height normalization.
\end{proof}

\begin{corollary}[The torsion-trace implication]\label{cm:torsion-trace}
In the CM elliptic setting, suppose only that the functional sign is
odd. Choose a split Heegner companion \(K'\) with nonzero central
value and form the same horizontal exponential coordinate \(Z\).
The prime choices need satisfy the spectral conditions above, but
need not satisfy the additional prescription involving a rational
point \(P\). If the projected Hilbert trace \(P_0\) is torsion,
then \(Z'=0\).
\end{corollary}

\begin{proof}
The spectral calculation \eqref{cm:eq:spectral-derivative} remains
valid: its proof used neither a simple zero of \(L(f_B,s)\) nor a
point generating the Selmer line.
For torsion \(P_0\), its tame height on the right is zero.
All factors multiplying \(Z'\) on the left are nonzero, including
\(L(f_B\otimes\varepsilon_{K'},1)\). Hence \(Z'=0\).
\end{proof} 


\subsection{Central formulas in known analytic rank at most one}

Let \(E/\mathbb Q\) be CM, and put
\(M=T_2(E^\vee)\otimes\mathcal O\).
Use the unsmoothed class \(z=z_E\) with fixed support as in
\eqref{eq:C2}--\eqref{eq:C4}.
Let \(a\) be a primitive invariant Betti homology vector and
\(a^*(1)\) its polarization image on dual Tate.

\begin{lemma}[The exact CM lattice index]\label{cm:lattice-index}
If \(E\) has analytic order zero or one, then
\[
 d\bigl(C^+(M),a^*(1)\wedge z\bigr)=v_2(\gamma(a)).
 \tag{C5}\label{eq:C5}
\]
\end{lemma}

\begin{proof}
Ordinary global cohomology of \(M\) (or \(M^*\), commensurable with it) rationally has a line
in degree one only. Indeed Kummer Selmer has the rank given by classical GZK; strict dual
localization has zero kernel inside it by nonzero local log in rank one, and local \(H^2\) at
finite places vanishes by dual invariant vanishing. Thus Poitou--Tate kills rational \(H^2\)
globally and the ranks follow by Euler characteristic. By \eqref{eq:C2} and Shapiro, a
determinant generator for the positive complex of \(M^*\), written in generic cohomology, is
\((\gamma-c\gamma)(1)\wedge z\), the first vector denoting real boundary and the second lifted
rationally. Indeed the real action on \(M^*\) is regular with a free invariant basis as
indicated, taking the twist sign into account.

We normalize \(v_2(2)=1\). To calculate the index, the positive Euler length of a finite
quotient in comparing these lattices is minus its module length. Filter it residually (use a
common sublattice or scalar multiple); constituents by CM semisimplification are trivial or
irreducible induced from one nontrivial residual line on \(L\), the conjugate line then being
its inverse. For a trivial constituent this Euler calculation is global squareclasses and
Brauer reciprocity: positive \(h^1,h^2\) are \(|S_f|, |S_f|-1\), respectively, including the
ordinary real restriction in defining positive, and there are no other terms. The
arithmetic-cohomology comparison for squareclasses/Brauer follows as in the diagram
conventions (the universal cover contains imaginary layers). For the induced line use Shapiro
and the ordinary Euler characteristic \(-1\) on \(L\), with one real invariant line further
removed. Additivity proves the length claim. The determinant change of module bases from \(M\)
to \(M^*\) has valuation \(v_2(\mu)\) by actual polarization unimodularity; thus the change of
\(d\) on fixed positive tensors from \(M^*\) to \(M\) is \(+v_2(\mu)\), by the length
calculation. And \(a^*=\pm(\gamma(a)/\mu)(\gamma-c\gamma)\) before twist, by evaluation
duality and anti-invariance. This proves \eqref{eq:C5}.
\end{proof}

\begin{corollary}[The CM leading coefficient in known low rank]
\label{cm:known-low-rank}
If \(E/\mathbb Q\) is CM and \(\an(E)\le1\), then \(X(E)=0\).
\end{corollary}

\begin{proof}
In rank zero \eqref{eq:C3} gives for \(\exp^*z\) the scalar \(e_S L(E,1)/p_\gamma\) relative
to \(\omega\). In rank one \eqref{eq:C4} gives, using any free Mordell--Weil basis \(P\),
\[
 v_2\left(\frac{\log_\omega(\lambda_{E,*}^{-1}z)}{(\log_\omega P)^2}\right)
   =v_2\left(e_S\frac{L'(E,1)}{p_\gamma H(P)}\right)
\]
in the period-normalized sense there. Since \(p_\gamma\gamma(a)=\pm\Omega_0\) (connected
period), these are exactly the log/Haar inputs of the central positive calculations giving
\eqref{eq:L5} and \eqref{eq:P1}, with a common multiplier of valuation \(v_2(\gamma(a))\) and
no smoothing. Explicitly the determinant raw valuation formula there uses the integral real
boundary basis \(a^*(1)\), subtracts \(\sum_{v\in S_f}\tau_v+\operatorname{length}H^1(\mathbb
R,T_2 E)\) and the dyadic \(\Sha\)-length, adds \(2\tau_g\), and inserts \(v_2\) of the above
scalar (exponential scalar at rank zero) plus the local log-image exponent \(b\). Haar absorbs
\(v_2(e_S)\) and \(b\), giving the Tamagawa factors and the real component correction to the
\emph{whole} period. Thus \(d(C^+(M),a^*(1)\wedge z)=X(E)+v_2(\gamma(a))\). This calculation
extends scalars by flatness, and if the generic elements require denominators to lift
integrally the length formula applies first to a multiple, then by homogeneity. Comparing to
\eqref{eq:C5} proves \(X(E)=0\). The basis \(P\) used here is a basis of the full
Mordell--Weil group modulo torsion; the calculation therefore uses
the regulator stipulated in Theorem~\ref{thm:main}.
\end{proof} 

\begin{corollary}[CM corank zero]\label{cm:corank-zero}
If \(E/\mathbb Q\) is CM and \(s_2(E)=0\), then
\(\an(E)=0\) and \(X(E)=0\).
\end{corollary}

\begin{proof}
Poitou--Tate gives rational concentration of ordinary global
cohomology in one degree-one line. The zero Selmer hypothesis makes
its map to the dyadic singular quotient injective, hence nonzero.
By \eqref{eq:C2}, \(z\) generates that line; by \eqref{eq:C3},
its nonzero singular localization implies \(L(E,1)\ne0\).
Apply Corollary~\ref{cm:known-low-rank}.
\end{proof}

\subsection{The paired CM determinant under a simple analytic zero}

\begin{proposition}[Paired CM comparison]\label{cm:paired}
Let \(A=e_fJ_0(N)\) have the normalization \eqref{eq:C1}. Suppose \(K\) has odd fundamental
discriminant \(-D<-4\), splitting \(2N\), and \(L(f/K,s)\) has a simple zero at 1. Use \(S\)
above finite rational primes splitting in \(K\), including \(2N\) in the prime-support sense,
archimedean places understood, and \(e_S\) for the omission multiplier of \(f\) over \(\mathbb
Q\). Put \(P=e_f P_X\) on \(A(K)_{\mathbb Q}\) for the Hilbert trace of \eqref{eq:GZ}. Over
\(K\) consider the complex \(C_*\) on \(M^*\). Fix one of the two primes \(w\mid2\),
impose the strict condition there, and impose full conditions at every other finite
prime of the support, including \(\bar w\). Then \(C_*\) is rationally acyclic and its characteristic determinant
has valuation
\[
                 -d(C_*,1)=2v_2(e_S\log_{\omega,w} P).               \tag{C-pair}\label{eq:C-pair}
\]
The differential in this identity is the normalized newform
differential, as in \eqref{eq:M3}.
\end{proposition}

\begin{proof}
Only the assumed simple zero of \(L(f/K,s)\), the Heegner rank
argument \eqref{eq:K0}, and the known-simple-zero comparison
of Proposition~\ref{cm:tame-log} are used here. In particular,
this proof does not use the CM corank-one converse proved below.

\par\smallskip\noindent\textbf{Concentration and the integral induction sequence.}\enspace

Here \(K\ne L\) and the plane remains absolutely irreducible over \(K\) as in \eqref{eq:C4}.
\eqref{eq:K0} gives ranks \(0,1\) on the even, odd functional-sign constituents over \(\mathbb
Q\), the latter generated on primal Tate by \(P\) in Kummer, viewed as a rational point up to
denominator of the sign-appropriate curve/variety by untwisting (\(-w(f)\) conjugation sign).
Logarithm on that generator is nonzero as before. Thus ordinary global cohomology for each
constituent over \(\mathbb Q\), on \(M^*\) or its quadratic twist by \(\varepsilon_K\), has
just a line in degree one rationally by the same Poitou--Tate test. Initially include also the
primes of \(D\), writing \(e,e'\) for the resulting omission multipliers of
\(f,f\otimes\varepsilon_K\). \eqref{eq:C2} gives generators \(z_0,z_1\) for these inverse
determinants by Shapiro from \(L\). The integral induction exact sequence with summands (as
sub and quotient) \(M^*,M^*\varepsilon_K\) on diagonal and difference, and Shapiro from \(K\),
then give
\(\operatorname{res} z_0\wedge\operatorname{res} z_1/2\) as ordinary inverse determinant
generator over \(K\). Indeed anti-diagonal insertion, corresponding to twisted restriction,
composes to twice the identity on the quotient.

\par\smallskip\noindent\textbf{The localized determinant.}\enspace By \eqref{eq:C3} and \eqref{eq:C4}, localization of
this pair at \(w\) is rationally an isomorphism on degree-one spaces, by nonzero exponential
on the rank-zero factor and finite nonzero Kummer log on the other; local cohomology likewise
has no other degrees rationally. Its inverse determinant volume is self-dual via the cup and
invariant using the alternating Tate form divided by \(\mu\), by integral local duality. Hence
the \(d\)-value of this localized wedge is the valuation of the cross pairing of the two
localized vectors (including division of the wedge by 2). Indeed the discriminant of the
determinant pairing on that basis is up to sign the square of the cross value since the finite
vector is isotropic. The triangle for strict thus gives before removing \(D\)
\[
   -d(C_{*,{\rm big}},1)
       =v_2\left(\frac{e e'\,(\log_{\omega,w}P)^2 L'(f/K,1)}
                              {2\mu p_\gamma^2 H_{\lambda_A,F}(P,P)}\right).
\]
The mixed cup before division by \(\mu\) evaluates the dual exponential against primal Kummer
log by \(\lambda_A^{-1}\) in the \(F\)-pairing. The sign-appropriate twist in \eqref{eq:C3} or
\eqref{eq:C4} is identified geometrically over \(K\) and at \(w\), with the indicated
\(\omega,p_\gamma\), and transports the \emph{absolute} pairing unchanged. The exact period
cancellation is
\[
 \frac{L'(f/K,1)}
      {2\mu p_\gamma^2H_{\lambda_A,F}(P,P)}
       =\frac{\pm i}{\sqrt D},
\]
by \eqref{eq:C1} and \eqref{eq:GZ}.
It is a dyadic unit because \(D\) is odd.

\par\smallskip\noindent\textbf{Removing the discriminant primes.}\enspace Removing the primes of \(D\) by localization
subtracts the valuations of just their \(f\)-Euler multipliers (good unramified representation
over each ramified place, singular Frobenius block on \(M^*(-1)\); the twist Euler factors
there are 1). The remaining primes contribute \(e_S^2\) since split, proving
\eqref{eq:C-pair}. This proves the claimed determinant formula.
\end{proof} 


\subsection{Horizontal divisors of the elliptic-unit class}

The remaining CM converse needs two kinds of information about an
elliptic-unit determinant coordinate. The individual-character theorem
controls its values at finite characters. The following lemma controls
its divisors away from two. Its hypothesis of generic concentration will
be established by the Bockstein argument in the next subsection.

\begin{lemma}[Horizontal nonnegativity for elliptic units]
\label{cm:horizontal-nonnegative}
Let \(E/\mathbb Q\) be CM by \(L\), with character lattice
\(T_\alpha\), and fix a support \(S_f\) containing two and all bad
primes. Choose good primes \(r_i\) split in \(L\), tending to one
dyadically, and even surjective residue characters
\[
 \lambda_i:(\mathbb Z/r_i\mathbb Z)^\times
                      \longrightarrow\mathbb Z/2^{m_i}\mathbb Z,
 \qquad m_i\longrightarrow\infty.
\]
Suppose these characters kill the fixed support and smoothing norms,
and the Tate Frobenii at \(r_i\) tend to a determinant-one element
with no root-of-unity eigenvalues. Use the compatible residue exponents
and finite-precision limits of the CM horizontal construction above;
no analytic-rank hypothesis is imposed here.

Put \(R=\mathcal O[[v]]\), and let \(C\) be the limiting ordinary
global complex over \(L\) on \(T_\alpha\), with scalar action
\((1+v)^{\lambda_i}\) or its consistent inverse, allowing the fixed
support and both places over \(r_i\). Let \(Y\) be the elliptic-unit
class with one fixed nonzero smoothing multiplier, including a fixed
integer clearing denominators when needed. Suppose \(C\) has one generic cohomology
line in degree one and no other generic cohomology, and that \(Y\)
is nonzero on this line. For an integral generator of \(\mathcal D(C)\),
write \(U\in\operatorname{Frac}(R)\) for the coordinate of \(Y\).
Then
\[
 \operatorname{ord}_{\mathfrak p}U\ge0
 \qquad\text{for every height-one }\mathfrak p\subset R
                    \text{ with }2\notin\mathfrak p.
\]
\end{lemma}

\begin{proof}
We work over the DVR \(R_{\mathfrak p}\), for a height-one prime
\(\mathfrak p\) not containing \(2\), and prove nonnegativity by
successive line switches over \(L\). On the fiber of \(R_{\mathfrak p}\), all local dual invariants at the original
allowed finite places vanish: at fixed places by constant local Tate torsion finiteness, at
old moving places by inertia if \(v\ne0\) and the Frobenius conditions if \(v=0\). Global
invariants also vanish by the kernel calculation below. Thus for the current primal problem
(initially full), with unramified or transverse conditions as below at auxiliary new places
and complementary conditions on the Tate-dual problem, we have
\(h=\dim H^1=1+h^*\)
on the fiber by Euler and Poitou--Tate, where \(h^*\) is the dual Selmer dimension in degree
one, injecting into global cohomology also (full local \(H^0\) included on the switch planes).
The Euler assertion persists under switches by the equal local quotient ranks; unramified
inflation recovers the old complexes.

If \(h>1\), we arrange a fresh place \(Q=Q_i\) outside the current support, splitting in stage
scalar fields and in preceding derivative fields, with \(\mathrm N Q\to1\), Tate Frobenius
tending to 1, and with both a primitive reduction of the generic primal cohomology line and
some nonzero dual fiber class evaluating nontrivially on its Frobenius. Over \(R_{\mathfrak
p}\) the local evaluations are a split finite axis (value on Frobenius) and a singular axis
(on a tame generator), of rank one each; conditions include degree zero and omit degree two.
Each pure axis pairs trivially with the same opposite-side axis, with perfect mixed pairing:
the limit action is trivial, one uses increasing roots-of-unity precision and the
tame/invariant computation of the diagram conventions. Thus these are the local comparisons of
Lemma~\ref{lem:foundations-line}, with the unramified plane the pure finite one.

\par\smallskip\noindent\textbf{Class-group corrections in the auxiliary ray fields.}\enspace
Denote by \(B_i/L\) the base cyclic extension from \(\lambda_i\), totally ramified at both
places over \(r_i\). Preceding switches will use cyclic \(F_Q/L\) of order \(s_Q=2^{A_i}\)
(parameters possibly different per \(Q\)), totally ramified just at \(Q\), unramified
elsewhere, with each auxiliary place split in the base and all the other derivative fields.
Choose fixed prime ideals \(I_j\) giving cyclic generators of the class group (cyclic factors
with orders \(h_j\)), \(I_j^{h_j}=(x_j)\). For old auxiliary \(Q\) choose \((y_Q)=Q\prod
I_j^{-n_j}\). Avoid these \(I_j\) in choosing new places. Include in \(D_i/L\) all \(2\)-power
roots of unity, dyadic radicals of \(x_j,y_Q\), and \(B_i\) and all old \(F_Q\). Splitting at
a new place to sufficiently high finite depth on this data permits \(F_{Q_{\rm new}}\) as
indicated of any required finite \(2^A\) order for that depth: take a primitive residue
exponent killing units and all \(x_j,y_Q\), prescribe it on principal ideals, and extend to
ideals away from the modulus by killing \(I_j\), which respects the class-power relations. It
kills each old \(Q\); ray reciprocity gives the assertion and total ramification from the
residue image. Roots for killing the finite global units can be included by the roots-of-unity
congruence.

\par\smallskip\noindent\textbf{Evaluation injectivity with uniform auxiliary data.}\enspace
In the product evaluation on \(\prod_i G_L\) for global cohomology classes in these problems,
restriction is injective to the exact joint kernel fixing all \(D_i\) and the constant Tate
actions termwise. In fact set \(K_T\) the extension by the character and cyclotomic Tate
actions, abelian over \(L\) (here \(T=T_\alpha\)). Its intersection with
\(D_i^0:=L(\mu_{2^\infty}) B_i\prod_Q F_Q\) is just \(L(\mu_{2^\infty})\), by unramifiedness
at these good odd places and separate total ramification of the cyclic factors. Indeed the
inertias surject separately to those whole factors alone, which are jointly independent by the
same ramification. The abelian part of \(D_i\) over \(L\) has bounded exponent over \(D_i^0\),
by one fixed nonidentity cyclotomic action on its radical translation subgroup over \(D_i^0\).
Consequently the Tate image over \(D_i\) still contains fixed powers (uniformly in \(i\)) of
the image over \(L(\mu_{2^\infty})\). To see this extend an automorphism of \(K_T\) over the
latter tower identically on \(D_i^0\) and then lift to the compositum with the abelian part;
the bounded power kills that part, hence kills \(D_i\cap K_T\). This gives a lift fixing
\(D_i\). The \(T\)-character on the cyclotomic kernel has infinite image by the CM character
argument. Thus some sequence trivial on the \(D_i\) acts by a fixed infinite-order scalar on
\(T\) (inverse on the dual), central in the product quotient by the joint kernel. This proves
injection by the central scalar test in characteristic zero, using also finite-list evaluation
injectivity of the diagram conventions, and proves global invariant vanishing. On that kernel
the values of the two indicated classes to be tested are additive, and neither test is
identically zero; take them simultaneously nonzero.

\par\smallskip\noindent\textbf{Chebotarev and precision margins.}\enspace
Keeping elements trivial on the \(D_i\), perturb the resulting choices by arbitrarily deep
powers of the preceding scalar lifts, so as to approach Tate identity with \(\rho_T-1\)
valuation tending to infinity but finite. Powers may preserve all finite old cochain
congruences to increasing precision. Take \(A_i\) much larger still. Chebotarev at finite
levels now gives \(Q_{\rm new}\) with the stated evaluations and splitting data including
residue prescriptions, with
\begin{equation}\label{cm:eq:precision}
\begin{gathered}
 v_2(\mathrm NQ_{\rm new}-1)>A_i,\qquad
 A_i-v_2(F_i-1)\longrightarrow\infty,\\
 v_2(F_i-1)\longrightarrow\infty,\qquad
 F_i=\rho_T(\mathrm{Fr}_{Q_{\rm new}}).
 \end{gathered}
\end{equation}
All fields and class prescriptions may use deeper finite precision in \(D_i\) as necessary
after choosing \(A_i\). The Frobenius used on old evaluations is unaffected by a local inertia
adjustment at the new place. This proves the required compatibility inductively; the new
conductor adds just one prime \emph{ideal} at a time in the norm over \(L\).

\par\smallskip\noindent\textbf{Derivative classes and bounded descent.}\enspace
For subsets \(I\) of places switched so far, use smoothed elliptic-unit classes evaluated
first in \(T\) over \(H_I=B_i\prod_{Q\in I}F_Q\), with support \(Q\in I\) added to the base
support and the corresponding elliptic-unit conductors. Apply to them \(\prod_{Q\in
I}\mathfrak D_Q\), \(\mathfrak D_Q=\sum_{j=1}^{s_Q-1} j\sigma_Q^j\), with \(\sigma_Q\) the
image of an oriented tame generator at \(Q\). At total splitting in all other fields the lower
norm on deleting \(Q\) uses \(p_Q=1-F_i/\mathrm N Q\) (with \(F_i\) at this \(Q\)), so
\(p_Q/(F_i-1)\to-1\). These are the elliptic-unit smoothing, twist and prime-ideal Euler
distributions specified at \eqref{eq:C2}–\eqref{eq:C3}. The group-ring identity
\((\sigma_Q-1)\mathfrak D_Q=s_Q-N_Q\) (\(N_Q\) the group sum) gives global invariance in
reduced cohomology to increasing precision, by the norm relation and the margins.

Uniformly on \(H_I\), invariants in \(T/2^n\) and \(H^1(T)\)-torsion are killed by a fixed
nonzero multiplier (disjoint Tate data by total ramification as above). Thus any
rational-versus-integral torsion errors in exact norm comparisons can first be killed in
integral cohomology, and another fixed multiplier makes the invariant reduced classes descend
to \(B_i\) by inflation-restriction (obstruction in cohomology of bounded-exponent
invariants). These choices are common across subsets; uniqueness loss is likewise bounded.

Weight-corestrict the descended classes from \(B_i\) by the horizontal character, modulo
increasing Artin precisions. Denote the resulting limit classes by \(Y_I\), with
\(Y_\varnothing\) the common multiple of \(Y\); they can be used successively with old choices
retained, slowing working precisions if needed.

\par\smallskip\noindent\textbf{The exact finite--singular interchange.}\enspace
For \(Q\in I\) these have the relations in the limit
\begin{equation}\label{cm:eq:interchange}
f_{\rm loc}(Y_I)=0,\qquad
                     s_{\rm loc}(Y_I)=\pm f_{\rm loc}(Y_{I-\{Q\}})
\end{equation}
in integral axes as above. Indeed before descent, with other derivatives applied and the
common multiplier understood, use integral cocycles \(b_Q\) upstairs and \(b\) below (over
\(H_{I-\{Q\}}\) and restricted up). They vanish on their respective local inertias above
\(Q\): the local coefficients are unramified, the residue cardinality is still \(\mathrm N
Q\), and \(F_i\ne\mathrm N Q\), excluding any nonzero inertia map to torsion-free \(T\). Use
\(\sigma_Q\) lifted in inertia, so \(\rho_T(\sigma_Q)=1\) and \(b_Q(\sigma_Q^{s_Q})=0\). Take
the residue lift \(\mathrm{Fr}\) trivial in all the stage finite cyclic fields here (adjust by
inertia in the totally ramified factor). The exact norm identity on upper global cochains has
\begin{equation}\label{cm:eq:cochain-norm}
N_Q b_Q=p_Q b+d e,\qquad
           (F_i-1)e=s_Q b_Q(\mathrm{Fr})-p_Q b(\mathrm{Fr})
\end{equation}
for some integral \(e\), by the cohomological distributions and unramified evaluations. We can
work separately with \(Q\) outermost in the derivative (orders commute on classes). Also
\((\sigma_Q-1)\mathfrak D_Q b_Q=s_Q b_Q-N_Q b_Q\) holds on cochains up there, since the extra
inner action by \(\sigma_Q^{s_Q}\) fixes this cocycle. Match the upper restriction of the
reduced descended cocycle with \(\mathfrak D_Q b_Q\) up to coboundary. Applying \(\sigma_Q-1\)
identifies the coboundary of its value at \(\sigma_Q\) with \(-d e\) modulo precision (the
mismatch coboundary is \(\sigma_Q\)-fixed). Invariant errors cost bounded precision, so its
value tends congruentially to \(-e\). The exact displayed equation, taken before reduction,
makes this agree up to sign in the limit with \(b(\mathrm{Fr})\) by the margins, which also
calculates the lower descent value there (a coboundary adjustment costs a term divisible by
\(F_i-1\)). On \(\mathrm{Fr}\) itself the upper derivative \(\mathfrak D_Q b_Q\) has value
\(s_Q(s_Q-1)b_Q(\mathrm{Fr})/2\), tending to zero, and descent adjustment again disappears.
Total splitting makes the same argument hold on all translates with their transported
generators, since the cyclic fields are abelian over \(L\). Thus the weighted traces retain
these identities. Unswitching conditions at other new places not in a given subset are
unramified for the corresponding classes by support. The cohomology classes with the indicated
local coordinates lift to the stated Selmer problems over the test DVR by the localization
triangles.

\par\smallskip\noindent\textbf{Reversing the switches.}\enspace
This is the unit line isomorphism on evaluations needed in the line switch, with classes
integral at the test DVR. By the nonzero primal and dual tests, each such switch drops \(h\)
by one on its fiber, preserving the generic line and its nonzero derivative by the triangles
from the lower condition. Iterate with the new primes as above. At \(h=1\) integrality gives
nonnegative order of the derivative coordinate there (one-line minimal complex). Reversing the
determinant equalities proves \(\operatorname{ord}_{\mathfrak p} U\ge0\), since the common
multiplier is a unit here.

\end{proof}

\subsection{The CM converse from Selmer corank one}

\begin{proposition}[CM corank one]\label{cm:corank-one}
If \(E/\mathbb Q\) is CM and \(s_2(E)=1\), then
\(\an(E)=1\) and \(X(E)=0\).
\end{proposition}

\begin{proof}
\par\smallskip\noindent\textbf{An abstract Selmer class and a transverse tame direction.}\enspace

Parity gives odd functional sign. We will construct a tame exponential
series \(Z(v)\) with \(Z(0)=0\) and prove \(Z'(0)\ne0\).
Corollary~\ref{cm:torsion-trace} will then make the Hilbert trace
nontorsion, so Gross--Zagier will detect the analytic simple zero.
The argument for \(Z'(0)\ne0\) has two parts: a nonzero Selmer
Bockstein gives a simple zero of the finite Selmer determinant, and
elliptic-unit comparisons make the ordinary determinant coordinate a
unit at the same characteristic-zero center.

Choose a nonvanishing split companion \(K'\) and the spectral
data of Corollary~\ref{cm:torsion-trace}.
At this stage no simple analytic zero or rational point spanning
the Selmer line is assumed. Use a single tame prime sequence \(r_i\) there with variable
\(v\), order \(2^{m_i}\to\infty\), action by \((1+v)^{\lambda_i}\) or its inverse. We impose
additionally \eqref{eq:P2}. More explicitly take \(x,\beta,k\) as in the unknown-center test
\eqref{eq:P2}–\eqref{eq:P3}, on primal and polarization-dual actual Tate over \(\mathbb Q\)
and scalar twist (1) respectively, unramified-data cochains outside fixed support \(S\),
integral up to fixed scalings. Here \(x\ne0\) in rational Selmer,
\(\beta=\operatorname{pol}(x)\), \(dk=-\beta\cup x\), which exists by Poitou--Tate and local
Kummer isotropy as there. Identify rationally the dual Tate with \(M^*\) after extending
coefficients. On the kernel fixing full Tate as well as \(F_0\), roots and prime radicals of
the \hyperref[cm:construction:transversality]{Frobenius and height construction}, the \(x\)-values span the plane: the argument there (central
homothety over \(F_0\), exclusion of factoring the remaining values through cyclotomically
conjugated translations, normality and absolute irreducibility) needs only nonzero cohomology
of \(x\).

Choose \(g\) fixing the required auxiliary data and with determinant one and nonroot
eigenvalues as there, and adjust by a joint-kernel commutator to achieve
\begin{equation}\label{cm:eq:transverse-selmer}
k(g)-\beta(g)g_T(g_T-1)^{-1}x(g)\ne0 .
\end{equation}
The adjustment works by the same cochain identity. Choose \(r_i,\lambda_i\) by approximation
to this \(g\), with all requirements of the CM horizontal sequence above, including splitting
in \(L,F_0\), increasing residue-root/radical precision and killing fixed support and
smoothing norms.

\par\smallskip\noindent\textbf{The ordinary and finite limiting complexes.}\enspace Use the perfect limiting ordinary
global complex \(C\) over \(R=\mathcal O[[v]]\) on \(T_\alpha\) over \(L\) with this scalar
action, allowing fixed support and both places above \(r_i\); equivalently by Shapiro use
\(M^*\) over \(\mathbb Q\). Include the smoothed elliptic-unit class \(Y\), whose exponential
coordinate at the split variable specialization local at 2 (that is, the scalar action there
is trivial for the whole variable) is the bounded-denominator series \(Z\) obtained by the
\hyperref[cm:construction:exponential]{exponential-coordinate construction}, applied to this sequence.
Smoothing is fixed nonzero as in that calculation. All comparisons here use the ordinary
finite free and cup diagrams specified earlier. At fixed places the variable is trivial; at
the two old moving places there is primitive inertia exponent, and at \(v=0\) the
determinant-one Frobenius prescription and norm limit 1 make both ordinary local complexes
rationally acyclic with unramified terms also acyclic. Thus unramified inflation there
recovers fixed-support global cohomology rationally at the center.

\par\smallskip\noindent\textbf{The order-one Bockstein.}\enspace Over the DVR \(R_{(v)}\) impose further at 2 the constant
Kummer line on the induced representation, giving \(C_F\). Local cohomology there is the
constant plane in degree one, with singular quotient measured by \(\exp^*\). At \(v=0\)
Poitou--Tate gives just cohomology lines in degrees 1 and 2 for \(C_F\), by corank one and
invariant vanishing (odd local terms acyclic). Its single minimal differential has order one.
To prove this, suppose instead that its order were at least two
(including the case of zero differential). Then \(\beta\) lifts on induced coefficients modulo
\(v^2\), with finite localization at 2. This gives the same closed-cochain test as
\eqref{eq:P3} over \(\mathbb Q\), now with no supplementary primes. Namely after a common
multiplier, stage representatives modulo \(v^2,2^{n_i}\) to slower cofinal precision have
constant term the fixed multiple of \(\beta\); their derivative \(z_1\) gives \(z_1\cup x\)
plus the corresponding multiple of \(\lambda_i\cup k\) closed. Here identify commensurable
lattices by maps with bounded denominators. The lift and adjustment are by the contractions
and constant-term coboundary lift as in \eqref{eq:P3}, retaining the dyadic localization and
cup with \(x\) on constant local models there. Thus its invariant tends to zero at 2; other
fixed places contribute zero after limit by bounded coboundary denominators (at infinity up to
bounded torsion), and outside support the cochain is unramified. At \(r_i\), trivializing
unramified \(x\) by bounded-denominator \(w_i=(g_{i,T}-1)^{-1}x(g_i)\) as there, \(g_i\)
arithmetic Frobenius, gives the prescribed nonzero invariant limit from
\(\lambda_i\cup(k-\beta\cup w_i)\), up to orientation and nonzero multiplier. This uses
exactly the primitive compatible residue exponent and Frobenius convergence, and contradicts
global reciprocity. This proves the order assertion, without assuming a global point spans.

\par\smallskip\noindent\textbf{Specialization at high-order characters.}\enspace
Let \(J\) be the constant dyadic singular quotient, a free line in
degree one. The order-one differential makes \(C_F\) generically
acyclic, so the triangle \(C_F\to C\to J\) identifies the generic
cohomology of \(C\) with this single line. Write
\(U\in\operatorname{Frac}R\) for the determinant coordinate of \(Y\),
using a generator of \(\mathcal D(C)\). For all fixed
roots of unity \(\zeta\) of sufficiently high 2-power order we have
\begin{equation}\label{cm:eq:high-character}
v_2 U(\zeta-1)=v_2\Delta ,
\end{equation}
with the fixed smoothing. Indeed nonzero pivots for the generic cohomology dimensions remain
nonzero at these evaluations by one-variable preparation. At such a fixed root the actual
character at stages is eventually allowed and the corresponding matrices converge with bounded
pivot denominators, so those stage complexes have the same rational concentration (Euler
characteristic \(-1\)). Their coordinates then satisfy the valuation test by \eqref{eq:C2}
separately, increasing dyadic powers of the modulus if needed, and use the same \(\Delta\)
since the smoothing norm is killed. Determinant base change and coefficient limits now give
the assertion, including generic nonvanishing of \(U\). This uses the one-character
determinant theorem as specified above, not an integral comparison across the whole tame
group.

\par\smallskip\noindent\textbf{The determinant coordinate at the center.}\enspace
The Bockstein calculation and \eqref{cm:eq:high-character} verify
respectively generic concentration and nonvanishing of \(Y\).
Lemma~\ref{cm:horizontal-nonnegative} therefore applies to this very
horizontal sequence. It gives nonnegative order for \(U\) at every
height-one prime not over two.

Here is the one-variable consequence of these two comparisons. Write,
by Weierstrass preparation,
\[
 U=\pi^a\varepsilon(v)\frac{P(v)}{Q(v)},
 \qquad \varepsilon\in R^\times,
\]
where \(P,Q\) are coprime distinguished polynomials and \(a\in\mathbb Z\).
Horizontal nonnegativity forces \(Q=1\). For roots of unity of
sufficiently high two-power order, the leading term of \(P\) gives
\[
 v_2 U(\zeta-1)
   =a\,v_2(\pi)+\deg(P)\,v_2(\zeta-1).
\]
Equation~\eqref{cm:eq:high-character} makes this expression constant
as \(v_2(\zeta-1)\) tends to zero. Hence \(\deg(P)=0\), and in
particular \(\operatorname{ord}_v U=0\).

\par\smallskip\noindent\textbf{Detection of the simple zero.}\enspace
Over \(R_{(v)}\) the triangle from \(C_F\) to \(C\) now gives
\begin{equation}\label{cm:eq:center-order}
0=-1+\operatorname{ord}_{v} Z,
\end{equation}
using the simple differential test and the constant singular line. Hence \(Z'\ne0\).
Corollary~\ref{cm:torsion-trace} now shows that the projected
Hilbert trace is nontorsion. The Gross--Zagier formula
\eqref{eq:GZ}, with the nonzero companion central value, gives
\(\an(E)=1\). Corollary~\ref{cm:known-low-rank} then gives
\(X(E)=0\).
\end{proof}

Corollary~\ref{cm:corank-zero} and Proposition~\ref{cm:corank-one}
prove the CM comparison, Proposition~\ref{prop:cm}.
The paired statement, Proposition~\ref{cm:paired}, remains available
under its separate simple-analytic-zero hypothesis.
